%% parts/p05/independence-catalog.tex --- Part 5
\chapter{The Independence Catalog}\label{sec:p05-independence-catalog}

\noindent\textit{Reader note.}  This separately authored catalog follows the
exact inherited G\"odel, Rosser, L\"ob, and second-incompleteness chapters.  It
adds combinatorial and model-theoretic examples without replacing their fixed
sentences, proof predicates, or proof statuses.  Primary sources are attached
to the individual claims below.

Each entry separates three facts: truth in the standard natural numbers,
formalizability in a base theory, and unprovability from that base under a
stated consistency or soundness assumption. Combinatorial size alone is not an
independence proof.

\section{Finite combinatorics beyond \(\mathsf{PA}\)}

\begin{hbdefinition}{def:p05-paris-harrington}{Paris--Harrington principle}
For positive integers \(k,r,n\), the Paris--Harrington assertion says that
there is \(N\) such that every \(r\)-coloring of the \(k\)-element subsets of
\(\{1,\ldots,N\}\) has a homogeneous set \(H\) with
\(|H|\ge n\) and \(|H|\ge\min H\). The second inequality is the largeness
condition and is necessary to distinguish the principle from ordinary finite
Ramsey theory. It stands as a first-order arithmetical statement family,
differs by this self-referential size bound, and is used in the independence
theorem \cite{ParisHarrington1977}.
\index{Paris--Harrington principle}
\end{hbdefinition}

\begin{hbtheorem}{thm:p05-paris-harrington}{cited=ParisHarrington1977}{Paris--Harrington independence}
The Paris--Harrington principle is true in the standard natural numbers but is
not provable in \(\mathsf{PA}\), assuming \(\mathsf{PA}\) is consistent. More
precise formulations calibrate the principle through reflection and ordinal
growth; the displayed statement does not claim independence from stronger
systems.
\end{hbtheorem}

\begin{hbmethod}{def:p05-ph-indicator}{The indicator route to unprovability}
For a nonstandard model \(\mathcal M\entails\PA\), an \emph{indicator for
\(\PA\)} is a definable numerical function \(Y(a,b)\) with the following use:
if \(a<b\) and \(Y(a,b)\) is nonstandard in \(\mathcal M\), the indicator
construction supplies a proper initial segment \(I\subsetneq M\), closed under
the arithmetic operations, such that \(a\in I\), \(b\notin I\), and
\(I\entails\PA\).  This name is needed to distinguish that cut-producing
function from the finite coloring itself.  It stands as a metamathematical
device on a nonstandard model, differs from an ordinal rank or a least Ramsey
number, and is used only in the following architecture.

Paris and Harrington define such an indicator from the sizes of relatively
large homogeneous sets.  Assuming that \(\PA\) proved the displayed principle,
their nonstandard construction would permit parameters below \(a\) while
placing the least witness for the corresponding internal finite-coloring
instance at or above \(b\).  The cut \(I\) would then satisfy \(\PA\) but fail
that instance: its parameters and coloring code lie in \(I\), while no witness
does.  Hence there is a model of \(\PA+\neg\mathrm{PH}\), so completeness
rules out \(\PA\derives\mathrm{PH}\).  Equivalently, the least-witness
functions extracted from suitable diagonal Paris--Harrington instances
eventually dominate every function whose totality \(\PA\) proves.  The
indicator and domination details are those of
\cite{ParisHarrington1977}; the broader finite-combinatorial method is
analyzed in \cite{KanamoriMcAloon1987}.
\end{hbmethod}

\begin{hbdefinition}{def:p05-full-aca}{\(\mathsf{ACA}_0\) and the full-induction theory \(\mathsf{ACA}\)}
In the two-sorted language of numbers and sets of numbers,
\(\mathsf{ACA}_0\) has the number-sort axioms of \(\mathsf Q\), set
extensionality, \(\Sigma^0_1\)-induction with number and set parameters, and
the comprehension scheme
\[
 \exists X\,\forall n\,(n\in X\leftrightarrow\varphi(n,\bar a,\bar A))
\]
for every arithmetical \(\varphi\), with \(X\) not free in \(\varphi\).
Its first-order consequences are exactly those of \(\PA\).
The un-subscripted \(\mathsf{ACA}\) is \(\mathsf{ACA}_0\) plus induction for
every formula of the two-sorted language, including formulas with set
quantifiers and set parameters.  The two names are necessary because that
full induction is used in the uniform Ramsey proof but is absent from the
standard reverse-mathematics subsystem.  Both stand as formal second-order
theories; their nonredundant difference is the induction scheme, and
\(\mathsf{ACA}\) is used here only for the following upper calibration.
\end{hbdefinition}

\begin{hbtheorem}{thm:p05-ph-upper-calibration}{cited=ParisHarrington1977}{Second-order upper calibration}
The full-induction theory \(\mathsf{ACA}\) proves the uniform infinite Ramsey
theorem for all finite dimensions and color numbers.  The usual infinite-to-
finite argument, with the relative-largeness condition imposed after choosing
the infinite homogeneous set, yields the Paris--Harrington principle.
Bare \(\mathsf{ACA}_0\) does not prove that arithmetical principle: it is
arithmetically conservative over \(\PA\), whereas
Theorem~\ref{thm:p05-paris-harrington} gives \(\PA\)-unprovability.  Thus the
register phrase “ACA context” means the full-induction theory just defined,
not \(\mathsf{ACA}_0\).
\end{hbtheorem}

\begin{hbdefinition}{def:p05-goodstein}{Goodstein sequence}
For \(b\geq2\), the \emph{hereditary base-\(b\) expression} of \(m\) is its
ordinary base-\(b\) expansion with every exponent recursively written in
hereditary base \(b\).  Its digits remain the numerals \(0,\ldots,b-1\).  Let
\(B_{b\to b+1}(m)\) be the value obtained by replacing every occurrence of
the base symbol \(b\) in that expression by \(b+1\).

For each starting value \(m\), define one sequence \((g^m_s)_{s\in\N}\) by
\[
 g^m_0=m,\qquad
 g^m_{s+1}=
 \begin{cases}
   0, & g^m_s=0,\\
   B_{s+2\to s+3}(g^m_s)-1, & g^m_s>0.
 \end{cases}
\]
Thus stage \(s\) uses hereditary base \(s+2\), beginning with base \(2\), and
zero remains zero.  The bound starting base and stage schedule are necessary
because a different first base gives a different numerical sequence.  The
term stands for this primitive-recursively specified family, differs from
ordinary base conversion by recursive exponent conversion and subtraction,
and is used in the ordinal descent below \cite{Goodstein1944,KirbyParis1982}.
\index{Goodstein sequence}
\end{hbdefinition}

\begin{hbdefinition}{def:p05-epsilon-zero-local}{The ordinal bound \(\varepsilon_0\)}
Every ordinal below \(\varepsilon_0\) has a unique Cantor normal form
\(\omega^{\alpha_0}c_0+\cdots+\omega^{\alpha_k}c_k\), with
\(\alpha_0>\cdots>\alpha_k\), positive finite coefficients, and exponents
again below \(\varepsilon_0\).  Here
\[
 \varepsilon_0=\min\{\alpha>0:\omega^\alpha=\alpha\}
 =\sup\{\omega,\omega^\omega,
          \omega^{\omega^\omega},\ldots\}.
\]
This local definition is needed because Part V uses the ordinal before the
later systematic ordinal chapter.  It stands as an ordinal, differs from the
finite hereditary-base expressions by replacing their base with \(\omega\),
and is used only to rank the two termination processes below.
\end{hbdefinition}

\begin{hbtheorem}{thm:p05-goodstein-termination}{cited=Goodstein1944}{Goodstein termination}
For every \(m\in\N\), some \(s\) satisfies \(g^m_s=0\).
\end{hbtheorem}

\begin{proofboundary}{def:p05-goodstein-descent-boundary}{Goodstein descent architecture}
For \(b\geq2\), let \(o_b(n)\) be the ordinal obtained from the hereditary
base-\(b\) expression of \(n\) by replacing \(b\) with \(\omega\).  Structural
induction on that expression gives
\[
 n>0\quad\Longrightarrow\quad
 o_{b+1}(B_{b\to b+1}(n)-1)<o_b(n)<\varepsilon_0.       \tag{G}
\]
Indeed, base change preserves the Cantor-term shape.  Subtracting one changes
the lowest nonzero term; after recursively borrowing through lower exponents,
the first Cantor-normal-form position at which the two ordinals differ has a
smaller exponent or coefficient on the left.  All newly introduced lower
terms are therefore dominated by the unchanged higher prefix.  Applying (G)
with \(b=s+2\) shows that every positive Goodstein step strictly lowers
\(o_{s+2}(g^m_s)\).  An infinite nonterminating sequence would give an
infinite descending sequence of ordinals, impossible for ordinals below
\(\varepsilon_0\).  This records the load-bearing rank and decrease.  A local
complete proof would additionally spell out every recursive borrowing case in
(G) and the well-foundedness argument for the chosen ordinal notation system;
those details remain carried by \cite{Goodstein1944,KirbyParis1982}, in
agreement with the theorem's cited status.
\end{proofboundary}

\begin{hbtheorem}{thm:p05-goodstein-independence}{cited=KirbyParis1982}{Goodstein independence}
The universal closure \(\forall m\,\exists s\,(g^m_s=0)\) is true by
Theorem~\ref{thm:p05-goodstein-termination}, but it is not provable in
\(\PA\) if \(\PA\) is consistent.
\end{hbtheorem}

\begin{hbdefinition}{def:p05-kirby-paris-hydra}{The Kirby--Paris hydra battle}
A \emph{hydra} is a finite rooted tree.  A head is a leaf together with its
incident edge.  At stage \(s\geq1\), Hercules chooses one head \(h\) with
parent \(p\) and deletes it.  If \(p\) is the root, nothing grows.  Otherwise
let \(q\) be the parent of \(p\).  In the decapitated tree, take the whole
subtree rooted at \(p\)---the part above the edge from \(q\) to \(p\)---and
attach to \(q\), in addition to the retained copy, exactly \(s\) disjoint new
copies of that subtree.  Hercules wins when only the root remains.

A strategy chooses the head from the current finite tree and stage number.
This exact tree, stage counter, and regrowth operation are necessary because
different hydra rules have different strength.  The definition stands as a
finite rewriting relation, differs from Goodstein's numerical base change,
and is used in the termination and unprovability theorem
\cite{KirbyParis1982}.
\index{hydra game!Kirby--Paris}
\end{hbdefinition}

\begin{hbtheorem}{thm:p05-hydra}{cited=KirbyParis1982}{Kirby--Paris hydra theorem}
Every battle governed by Definition~\ref{def:p05-kirby-paris-hydra} ends,
regardless of the chosen heads.  After coding finite trees and restricting
strategies to recursive functions, the first-order statement “every recursive
strategy wins” is not provable in \(\PA\), assuming \(\PA\) is consistent.
\end{hbtheorem}

\begin{proofboundary}{def:p05-hydra-boundary}{Hydra rank and independence boundary}
Assign ordinal \(0\) to every leaf.  If a nonleaf has children whose assigned
ordinals, in nonincreasing order, are
\(\alpha_1,\ldots,\alpha_k\), assign it
\(\omega^{\alpha_1}+\cdots+\omega^{\alpha_k}\); the root value is below
\(\varepsilon_0\).  A cut lowers the value at \(p\).  At \(q\), any finite
number of copies of the new term is still below the former
\(\omega^{\alpha}\)-term, and the strict decrease propagates to the root.
Therefore no infinite battle exists.  Kirby and Paris also define a recursive
strategy whose successive ranks realize the relevant fundamental-sequence
descent; proving that this strategy always terminates would prove a totality
principle beyond \(\PA\), yielding the cited unprovability result.  The full
node-by-node rank induction and the construction/calibration of that recursive
strategy remain in \cite{KirbyParis1982}; this paragraph fixes their shape but
does not claim a second local proof.
\end{proofboundary}

\section{Nonstandard models and computability}

The exact statement of Tennenbaum's theorem is
Theorem~\ref{thm:p05-tennenbaum}: no countable nonstandard model of
\(\mathsf{PA}\) has both addition and multiplication computable under a
presentation on \(\mathbb N\). The simultaneous-operation qualification does
not rule out computable presentations of selected reducts.

\begin{hbmethod}{def:p05-tennenbaum-method}{Tennenbaum's inseparability argument}
Choose disjoint c.e. sets \(A,B\subseteq\N\) with no computable separator, and
fix primitive-recursive enumerations of them.  This pair is needed because a
computable presentation would turn one internally finite code into a forbidden
external separator; it stands as external computability data, not as a set
inside the nonstandard model.

Suppose that \(\mathcal M\entails\PA\) is nonstandard and that both transported
operations on its domain \(\N\) are computable.  For each standard stage \(n\),
the two enumerations have no common value below stage \(n\).  Overspill gives a
nonstandard stage \(c\in M\) through which \(\mathcal M\) still sees them as
disjoint.  Using the internal recursion available in \(\PA\), form the product
\(a\) of the prime \(p_i\) for every value \(i\) enumerated into \(A\) before
stage \(c\).  For each standard \(i\),
\[
 i\in A\Longrightarrow \mathcal M\entails p_i\mid a,
 \qquad
 i\in B\Longrightarrow \mathcal M\entails p_i\nmid a.                 \tag{T}
\]
The second implication uses the internal disjointness through \(c\) and the
elementary prime-divisibility facts proved in \(\PA\).

It remains to verify the computability step rather than assume it.  From the
computable operations one computes every standard numeral and standard prime
in the presentation.  Given \(i\), search simultaneously over \(q\in\N\) and
the finitely many standard residues \(r<p_i\) until
\(a=p_i\times q+r\).  \(\PA\) proves division with remainder by each fixed
standard positive numeral, so exactly one such search succeeds; the answer
\(r=0\) decides whether \(p_i\mid a\).  Consequently
\[
 C=\{i\in\N:\mathcal M\entails p_i\mid a\}
\]
is computable.  By (T), \(A\subseteq C\) and \(B\cap C=\varnothing\), contrary
to computable inseparability.  Thus both operations cannot be computable in
one nonstandard presentation \cite{Tennenbaum1959,Kaye1991}.
\end{hbmethod}

\section{Consistency statements and presentation dependence}

\begin{hbdefinition}{def:p05-consistency-presentation}{Consistency statement for a presentation}
For a recursively axiomatized theory \(T\) with a chosen proof predicate
\(\mathsf{Proof}_T(p,q)\), put
\[
 \mathsf{Con}_{\mathsf{Proof}_T}(T)
 \defeq \neg\exists p\,\mathsf{Proof}_T
   (p,\overline{\ulcorner\bottom_A\urcorner}).
\]
The subscript is necessary because extensionally equivalent axiom sets or proof
predicates can yield arithmetically nonequivalent internal sentences in weak
bases. It stands as an arithmetical sentence tied to a presentation, differs
from the external assertion that no contradiction is derivable, and is used
to state intensionality warnings \cite{Feferman1960}. Whenever the chosen
\(\mathsf{Prov}_T\) is the existential closure of this particular
\(\mathsf{Proof}_T\), the displayed sentence is exactly
\(\mathsf{Con}(T)\) from
Definition~\ref{def:formal-consistency-sentence}. In the standard \(\PA\)
applications fixed in \cref{sec:second-incompleteness-01}, the proof formula
is specifically \(\mathsf{Prf}^{19}_{\PA}\), and consistency is
\(\neg\mathsf{Prv}^{19}_{\PA}
(\overline{\ulcorner\bottom_A\urcorner})\).
An earlier certified representative retains its own subscripted consistency
sentence; no internal equivalence with the selected direct presentation is
assumed. The subscript exposes this formula choice.
\index{consistency statement!presentation dependence}
\end{hbdefinition}

G\"odel's second incompleteness theorem applies to standard provability
predicates satisfying the derivability conditions. It does not license an
unqualified phrase ``\(T\) cannot prove its consistency'' for every formula
extensionally expressing consistency in the standard model. Feferman's work
on arithmetization and reflection shows why the numbering, axiom predicate,
base theory, and derivability conditions belong in the theorem statement
\cite{Feferman1960}.

\begingroup\small
\begin{longtable}{@{}p{0.20\textwidth}p{0.22\textwidth}p{0.25\textwidth}p{0.23\textwidth}@{}}
\caption{Arithmetic independence catalog.  “Independent”
means that neither the sentence nor its negation is provable only where both
directions are stated in the last column.}
\label{tab:p05-independence-catalog}\\
\toprule
Sentence or principle & Formal base & Standard truth or termination method &
Exact formal boundary \\
\midrule
\endfirsthead
\toprule
Sentence or principle & Formal base & Standard truth or termination method &
Exact formal boundary \\
\midrule
\endhead
G\"odel sentence \(G_T\) & c.e. \(T\supseteq\mathsf Q\) &
Fixed point for \(\neg\mathsf{Prov}_T\); standard truth follows from
consistency. &
\(T\notderives G_T\) under consistency; for
\(T\notderives\neg G_T\), use \(\Sigma^0_1\)-soundness or
\(\omega\)-consistency. \\
\addlinespace
Rosser sentence \(R_T\) & c.e. \(T\supseteq\mathsf Q\) &
Fixed point with bounded comparison of proof witnesses. &
Both \(T\notderives R_T\) and \(T\notderives\neg R_T\) follow from
consistency alone. \\
\addlinespace
\(\mathsf{Con}(T)\) & c.e. \(T\supseteq\mathsf Q\) with D1--D3 &
External consistency is the absence of a standard proof of
\(\bottom_A\). &
If \(T\) is consistent, \(T\notderives\mathsf{Con}(T)\); no two-sided
independence claim is made. \\
\addlinespace
Paris--Harrington & \(\PA\) &
Infinite Ramsey theory; indicator and nonstandard-cut analysis. &
True in \(\mathcal N\), unprovable in \(\PA\) under consistency; provable in
the full-induction \(\mathsf{ACA}\), not bare \(\mathsf{ACA}_0\). \\
\addlinespace
Goodstein termination & \(\PA\) &
Strict ordinal descent below \(\varepsilon_0\). &
The universal termination sentence is unprovable in \(\PA\) under
consistency. \\
\addlinespace
Kirby--Paris hydra termination & \(\PA\) &
Rule-specific tree rank below \(\varepsilon_0\). &
Every battle ends externally; “every recursive strategy wins” is unprovable
in \(\PA\) under consistency. \\
\bottomrule
\end{longtable}
\endgroup
